\scanpage{32}
\deleted{En d'autres termes, le \unsure{triangle}}
\deleted{\[H^i(W_0)=0\quad\text{si }i\ne\unsure{i|1}\]}
\deleted{\[\begin{gathered}\unsure{i=n+1}\\\unsure{2n-i=n-1}\end{gathered}\]}
\margin{\uline{\unsure{intéressant}}}
\[\left\{\begin{array}{ll}
H^i(W_0)\xrightarrow{\sim}H^{2n-i-1}(M)^\vee&\text{si }i\ge n+\unsure{1|2}\\
\boxed{H^i(M)\xrightarrow{\sim}H^i(W_0)}&\text{si }i\le n-2
\end{array}\right.\]
\margin{(épim si $i=\unsure{n-1|n}$)}
\margin{(donc $2n-i-1\le n-3$)}
\margin{(mono si $i=n-1$)}
Donc $H^{\unsure{q|*}}(M)$ est connu en termes de $H^{\unsure{q|*}}(W_0)$ sauf aux dimensions \deleted{$\unsure{n}$} $n-1,n$, \deleted{$\unsure{n}$}, où on a la suite exacte
\diagram{\begin{tikzcd}[column sep=small,row sep=large]
\deleted{0}\arrow[r]&\deleted{H^{n-2}(M)}\arrow[r]&\deleted{H^{n-2}(W_0)}\arrow[r]&\deleted{H^{\unsure{2n+1|n+1}}(M)^\vee}\inserted{0}\arrow[r]&H^{n-1}(M)\arrow[r]&H^{n-1}(W_0)\arrow[r,"\mathrm{surj}"]\arrow[d,no head,"\deleted{\unsure{\simeq}}"]&H^{\unsure{n}}(M)^\vee\arrow[d]\\
\deleted{0}&\deleted{H^{n-2}(M)^\vee}\arrow[l]&\deleted{H^{n+1}(W_0)}\arrow[l]&\deleted{H^{n+1}(M)}\inserted{0}\arrow[l]&H^{n-1}(M)^\vee\arrow[l]&H^n(W_0)\arrow[l]&H^n(M)\arrow[l,"\deleted{\unsure{\rho}}"]
\end{tikzcd}}

Supposons maintenant que $W$ \unsure{est} la sphère topologique. Alors on a vu que $H^i(W_0)=0$ si $i\ne0,n-1,n,2n-1$.

\uuline{Si donc} \textcircled{$n\ge3$}, on trouve aussi \uline{$H^{\inserted{i}\deleted{\unsure{n-2}}}(M)=0$}, \deleted{$\unsure{H^{n-2}(M)=0}$}
\inserted{si $i\ne0,n-1,n$, \deleted{$\unsure{2n}$}}

\deleted{et la suite exacte précédente devient}
\deleted{\diagram{\begin{tikzcd}[column sep=small,row sep=large]
&+&-&+&-\\
0\arrow[r]&H^{\unsure{n-1|n+1}}(M)^\vee\arrow[r,"u"]&H^{n-1}(M)\arrow[r]&H^{n-1}(W_0)\arrow[r]&H^n(M)^\vee\arrow[d,very thick]\\
0&H^{n+1}(M)\arrow[l]&H^{n-1}(M)^\vee\arrow[l,"u'"']&H^n(W_0)\arrow[l]&H^n(M)\arrow[l]\\
&-&+&-&+\\
&\unsure{H}
\end{tikzcd}}}
\mspage{31}
